세종이는 한양에서 유명한 트리 장인이다. 어떤 그래프라도 들고 오면 멋진 트리\footnote{트리는 $N$개의 정점과 $N-1$개의 간선으로 이루어진 무방향 연결 그래프이다.}로 탈바꿈해준다!

정점이 $N$개, 간선이 $M$개인 단순 그래프\footnote{단순 그래프는 각 간선이 서로 다른 두 정점을 이으며, 어떤 두 정점 $u$, $v$에 대해서도 $u$와 $v$를 잇는 간선이 둘 이상 존재하지 않는 그래프이다.}가 주어진다.

세종이는 정점 $1, 2, \ldots, N$ 중 두 정점을 잇는 $0$개 이상의 새로운 간선들을 적절히 추가하여 이 그래프를 트리로 만들 것이다.

세종이는 작업에 들어가기 전에 위 조건대로 그래프를 트리로 만드는 방법의 수를 가늠해 보려 한다. 세종이를 도와 그 방법의 가짓수가 $K$가지를 넘는지, 넘지 않는다면 정확한 가짓수까지 구해보자! 단, 추가한 간선 집합이 다를 경우에만 다른 방법으로 간주하며, 모든 간선에는 방향성이 없다.